\section{Introduction}
%Le ramassage des clients est une tâche essentielle pour les entreprises de transport, notamment pour les motos-taxis. En effet, il s'agit de garantir la satisfaction des clients en les assurant d'un transport rapide et fiable.
%Cependant, le ramassage des clients peut être un défi, car il nécessite de gérer de manière efficace une multitude de paramètres, tels que la position des clients, la disponibilité des motos, et les délais de livraison.
%Dans ce contexte, le projet moto que nous proposons vise à développer un algorithme de ramassage des clients qui respecte une certaine politique, tout en imitant le fonctionnement d'une moto dans la vie réelle.
%Quelle politique de ramassage des clients est la plus adaptée ?  Comment s'assurer que l'algorithme respecte la politique choisie ? Comment modéliser le fonctionnement d'une moto de manière à ce que l'algorithme soit robuste ? repondre a ces questions dans la suite permettra d’expliciter les détails de ce projet tout en defissant une politique et un mode de fonctionnement adéquat pour notre moto.




The transportation of customers is an essential activity for transportation companies, especially for motorcycle taxi providers. It is indeed about guaranteeing customer satisfaction by offering them fast and reliable transportation.
However, customer transportation poses a challenge, as it requires effectively managing a multitude of parameters, such as customer location, motorcycle availability, and delivery times.
In this context, the motorcycle project we are carrying out aims to develop a customer transportation algorithm that respects a certain policy, while simulating the natural operation of a motorcycle in reality.
The customer transportation policy consists of defining the selection criteria for customers to be transported, as well as the order of priority of requests. The customer transportation algorithm must be able to respect the defined policy, while optimizing the use of available resources. The operation of a motorcycle in reality includes the behavior and characteristics of a real motorcycle, such as: looking for a customer when the motorcycle is empty, going to a given destination when the motorcycle has a customer...etc. Thus, we ask ourselves the question of what problem does this problem situation raise? What description can we make of this problem? What modeling will best simulate reality? and how to implement this modeling to have the most optimal solution?
Answering these questions in the following will clarify the terms of this project while defining a policy and an appropriate mode of operation for our motorcycle.


%{lien vers le dépôt GITHUB}

\begin{multicols}{}
    \begin{itemize}
   
        \item \bgtext\href{https://github.com/Nameless0l/PROJET-TAXI.git}{https://github.com/Nameless0l/PROJET-TAXI.git}
    \end{itemize}
    \end{multicols}